\documentclass[a4paper,12pt]{article}
\usepackage{color}
\usepackage{graphicx, gensymb}
\usepackage{amsfonts, cancel, amssymb}
\usepackage{amsmath, amsthm, array, esvect, esint}
\usepackage{typearea}
\usepackage[a4paper,left=2cm,right=2cm,top=2cm,bottom=3cm,bindingoffset=5mm]{geometry}
\newcommand\numberthis{\addtocounter{equation}{1}\tag{\theequation}}
\newcommand{\bra}{}
\newcommand{\kom}{}
\newcommand{\brak}{}
\newcommand{\braket}{}
\newcommand{\intx}{}
\newcommand{\intr}{}
\newcommand{\kla}{}
\newcommand{\klam}{}
\newcommand{\klammer}{}
\newcommand{\oper}{}
\newcommand{\tecxt}{}
\newcommand{\Bra}{}
\newcommand{\Ket}{}

\begin{document}

\title{09 Elektronische Bandstruktur}
\author{Joachim Schwardt}
\date{\today}
\maketitle

\section*{Aufgabe 9.1: Sommerfeld-Entwicklung}
Sei $f(-x)=f(x)$, $t(-x)=t(x)$ und $g(-x)=-g(x)$.\\ 
Dann gilt für $a>0$ mit $\frac{1}{1+t^g}+\frac{1}{1+t^{-g}}=\frac{1+t^{-g}+1+t^g}{(1+t^g)(1+t^{-g})}=1$ folgende Rechenregel:
\begin{align*}
I&:=\intx\int_{-a}^{a}\underbrace{\frac{f}{1+t^g}}_{=:F(x)}\,\mathrm{d}x=\frac{1}{2}\intx\int_{-a}^{a}F(x)+F(-x)\,\mathrm{d}x+\frac{1}{2}\cancelto{0}{\intx\int_{-a}^{a}F(x)-F(-x)\,\mathrm{d}x}=\\
&=\frac{1}{2}\intx\int_{-a}^{a}f(x)\klam\left[ {\frac{1}{1+t(x)^{g(x)}}+\frac{1}{1+t(x)^{-g(x)}}} \right]\,\mathrm{d}x=\intx\int_{0}^{a}f(x)\,\mathrm{d}x
\end{align*}
Es ist $f(\epsilon)=\frac{1}{1+\exp\kla\left( {\frac{\epsilon -\mu (T)}{k_BT}} \right)}$ die Fermi-Funktion, $D(\epsilon)=\frac{m}{\pi^2\hbar^2}\sqrt{2m\epsilon}$ die Zustandsdichte und $a(\epsilon)$ beliebig. Betrachte das folgende Integral für kleine Temperaturen $T$ mit der Substitution $x:=\frac{\epsilon -\mu }{k_BT}$ und nutze $\frac{1}{1+e^{-x}}=1-\frac{1}{1+e^x}$:
\begin{align*}
A(T)&:=\intr\int_{\mathbb{R}^{ }}{H(\epsilon )f(\epsilon )}\,\mathrm{d}^{ }\epsilon = k_BT\intr\int_{\mathbb{R}^{ }}{\frac{H(\mu +k_BTx)}{1+e^x}}\,\mathrm{d}^{ }x=\\
&=k_BT\intx\int_{0}^{\infty}\frac{H(\mu +k_BTx)}{1+e^x}\,\mathrm{d}x+k_BT\intx\int_{-(-\infty)}^{-0}\frac{H(\mu -k_BTx)}{1+e^{-x}}\,(-\mathrm{d}x)=\\
&=k_BT\intx\int_{0}^{\infty}H(\mu -k_BTx)\,\mathrm{d}x+ k_BT\intx\int_{0}^{\infty}\frac{H(\mu +k_BTx)-H(\mu -k_BTx)}{1+e^x}\,\mathrm{d}x=\\
&=k_BT\intx\int_{\mu}^{-\infty }H(\epsilon)\,\frac{-\mathrm{d}\epsilon}{k_BT} + k_BT\intx\int_{0}^{\infty}\frac{2k_BTxH'(\mu)+\mathcal{O}[(k_BTx)^2]}{1+e^x}\,\mathrm{d}x\simeq\\
&\simeq\intx\int_{-\infty}^{\mu }H(\epsilon)\,\mathrm{d}\epsilon +2(k_BT)^2H'(\mu)\intx\int_{0}^{\infty}\frac{x}{1+e^x}\,\mathrm{d}x=\intx\int_{-\infty}^{\mu }H(\epsilon)\,\mathrm{d}\epsilon +\frac{\pi^2}{6} (k_BT)^2H'(\mu)
\end{align*}
Damit lässt sich auch folgendes Integral annähern:
\begin{align*}
B(T)&:=\intx\int_{0}^{\infty}H(\epsilon)(-\partial_\epsilon f)\,\mathrm{d}\epsilon =\intx\int_{0}^{\infty}f\partial_\epsilon H\,\mathrm{d}\epsilon\simeq H(\mu) +\frac{(\pi k_BT)^2}{6}\partial_\epsilon^2H(\mu)
\end{align*}
Sei nun $H(\epsilon)=D(\epsilon)$ und $A(T)=n=\frac{N}{V}$. Man findet $\alpha = \frac{\sqrt{8m^3}}{3\pi^2\hbar^2}$ und $\beta = \frac{\pi^2k_B^2}{8}$:
\begin{align*}
n&\simeq\intx\int_{0}^{\mu}\frac{m}{\pi^2\hbar^2}\sqrt{2m\epsilon}\,\mathrm{d}\epsilon +\frac{(\pi k_BT)^2}{6}\frac{m}{\pi^2\hbar^2}\partial_\epsilon\sqrt{2m\epsilon}\bigg\vert_\mu=\\
&=\frac{\sqrt{2m^3}}{\pi^2\hbar^2}\klam\left[ {\frac{2}{3}\sqrt{\mu^3}+\frac{(\pi k_BT)^2}{12\sqrt{\mu}}} \right]\equiv\alpha \mu^{3/2}+\alpha\beta T^2\frac{1}{\mu^{1/2}}
\end{align*}Es gilt $\mu_0\equiv \mu (T=0)=(n/\alpha)^{2/3}$. Umstellen der Gleichung und entwickeln in $T$ liefert dann:
\begin{align*}
\mu^{3/2}&=\frac{n}{\alpha}-\frac{\beta T^2}{\mu^{1/2}}=\mu_0^{3/2}\klam\left[ {1-\frac{\beta T^2}{\mu^{1/2}\mu_0^{3/2}}} \right]\\
\Rightarrow\ \ \ \mu&=\mu_0\klam\left[ {1-\frac{\beta T^2}{\mu^{1/2}\mu_0^{3/2}}} \right]^{2/3}\approx\mu_0\klam\left[ {1-\frac{2}{3}\frac{\beta T^2}{\mu_0^2 }} \right]=\mu_0\klam\left[ {1-\frac{1}{12}\kla\left( {\frac{\pi k_BT}{\mu_0}} \right)^2} \right]
\end{align*}
\section*{Aufgabe 9.2: Fermi-Gase in $d$ Dimensionen}
Es gilt mit $Z(k)=2\cdot\frac{V}{(2\pi)^d}$ folgender Zusammenhang zur Zustandsdichte (Faktor 2 kommt von Spin-Entartung), wobei aus der Dispersion $E(k)=\frac{\hbar^2k^2}{2m}$ noch $\tecxt\text{d}E=\frac{\hbar^2k}{m}\,\tecxt\text{d}k$ folgt:
\begin{align*}
Z(k)\,\tecxt\text{d}^dk&=2\cdot\frac{V}{(2\pi)^d}\,\frac{2\pi^{d/2}}{\Gamma\kla\left( {\frac{d}{2}} \right)}\cdot k^{d-1}\,\tecxt\text{d}k\overset{!}{=}D(E)\,\tecxt\text{d}E\\
\Rightarrow\ \ \ D(E)&=\frac{V\kla\left( {\frac{2mE}{\hbar^2}} \right)^{\frac{d-2}{2}}}{2^{d-2}\pi^{d/2}\Gamma\kla\left( {\frac{d}{2}} \right)}\frac{m}{\hbar^2}=\frac{Vm}{\hbar^2}\frac{1}{\pi^{d/2}\Gamma\kla\left( {\frac{d}{2}} \right)}\kla\left( {\frac{mE}{2\hbar^2}} \right)^{\frac{d}{2}-1}
\end{align*}
Der Fermi-Wellenvektor lässt sich aus der Bedingung $N=Z(k_F)\tecxt\text{Vol}_d\,(k_F)$ bestimmen:
\begin{align*}
N&\overset{!}{=}2\cdot\frac{V}{(2\pi)^d}\frac{\pi^{d/2}k_F^d}{\Gamma\kla\left( {1+\frac{d}{2}} \right)}& \Rightarrow && k_F&=2\sqrt{\pi}\klam\left[ {\frac{N \Gamma\kla\left( {1+\frac{d}{2}} \right)}{2V}} \right]^{1/d}
\end{align*}
Daraus folgen weitere Größen:
\begin{align*}
v_F&=\frac{\hbar k_F}{m}=\frac{2\hbar\sqrt{\pi}}{m}\klam\left[ {\frac{N \Gamma\kla\left( {1+\frac{d}{2}} \right)}{2V}} \right]^{1/d}\\
E_F&=\frac{\hbar^2k_F^2}{2m}=\frac{2\hbar^2\pi}{m}\klam\left[ {\frac{N \Gamma\kla\left( {1+\frac{d}{2}} \right)}{2V}} \right]^{2/d}
\end{align*}
Für die Zustandsdichte an der Fermi-Kante folgt dann:
\begin{align*}
D(E_F)&=\frac{Vm}{\hbar^2}\frac{1}{\pi^{d/2}\Gamma\kla\left( {\frac{d}{2}} \right)}\pi^{\frac{d}{2}-1}\klam\left[ {\frac{N\Gamma\kla\left( {1+\frac{d}{2}} \right)}{2V}} \right]^{1-\frac{2}{d}}=\frac{Nmd}{4\pi \hbar^2}\klam\left[ {\frac{2V}{N\Gamma\kla\left( {1+\frac{d}{2}} \right)}} \right]^{\frac{2}{d}}
\end{align*}
Für $E(k)=\hbar kv_F$ ändert sich die Zustandsdichte:
\begin{align*}
D(E)&=\frac{V}{2^{d-2}\pi^{d/2}\Gamma\kla\left( {\frac{d}{2}} \right)}\frac{1}{\hbar v_F}\kla\left( {\frac{E}{\hbar v_F}} \right)^{d-1}
\end{align*}
Der Fermi-Wellenvektor und damit auch die Fermi-Geschwindigkeit sind unabhängig von der Dispersion. Daraus folgt dann die Zustandsdichte an der Fermi-Kante:
\begin{align*}
D(E_F)&=\frac{V}{\hbar v_F2^{d-2}\pi^{d/2}\Gamma\kla\left( {\frac{d}{2}} \right)}2^{d-1}\pi^{\frac{d-1}{2}}\klam\left[ {\frac{N\Gamma\kla\left( {1+\frac{d}{2}} \right)}{2V}} \right]^{1-\frac{1}{d}}=\\
&=\frac{Nd}{2\hbar v_F\sqrt{\pi}}\klam\left[ {\frac{2V}{N\Gamma\kla\left( {1+\frac{d}{2}} \right)}} \right]^{1/d}=\frac{Nmd}{4\pi\hbar^2}\klam\left[ {\frac{2V}{N\Gamma\kla\left( {1+\frac{d}{2}} \right)}} \right]^{2/d}
\end{align*}
Allgemein gelten folgende Formeln:
\begin{align*}
N&=\intr\int_{\mathbb{R}^{d}}{Z(k)}\,\mathrm{d}^{d}k\kom\overset{\substack{{Z=\tecxt\text{const}} \\ |}}{=}Z\tecxt\text{Vol}_d\,(k)\\
D(E)&=\frac{\tecxt\text{d}N}{\tecxt\text{d}E}\\
N&=\intx\int_{0}^{E_F}D(E)\,\mathrm{d}E
\end{align*}
\section*{Aufgabe 9.3: Zweidimensionales Fermi-Gas}
Die mittlere Gesamtenergie berechnet sich unter Verwendung der Bedingung $N=Z(k_F)\tecxt\text{Vol}_d\,(k_F)$ zu
\begin{align*}
\bra\langle {E}\rangle&=\frac{E_\tecxt\text{ges}}{N}=\frac{1}{N}\sum_{k<k_F}\frac{\hbar^2k^2}{2m}\simeq\frac{1}{N}\intr\int_{\mathbb{R}^{d}}{Z(k)\frac{\hbar^2k^2}{2m}}\,\mathrm{d}^{d}k=\\
&=\frac{1}{\tecxt\text{Vol}_d(k_F)}\intx\int_{0}^{k_F}\frac{\hbar^2k^2}{2m}\,d\tecxt\text{Vol}_d(1)k^{d-1}\mathrm{d}k=\frac{d}{d+2}\frac{\hbar^2k_F^2}{2m}=\frac{d}{d+2}E_F
\end{align*}
Verwende noch $N=Z(k_F)\tecxt\text{Vol}_d\,(k_F)$
Einige Werte sind in der Tabelle aufgelistet.
\begin{table}[h!]
\centering
\begin{tabular}{c|c|c|c|c|c|c|c|c}
$d$&1&2&3&4&5&6&\dots &$\infty$\\ \hline
$\bra\langle {E}\rangle\,/\,E_F$&$\frac{1}{3}$&$\frac{1}{2}$&$\frac{3}{5}$&$\frac{2}{3}$&$\frac{5}{7}$&$\frac{3}{4}$&\dots &1\\
\end{tabular}
\end{table}
Der Druck kann über partielles Ableiten nach dem Volumen bestimmt werden. Oben wurde bereits der Zusammenhang $U\propto E_F\propto V^{-2/d}$ gefunden:
\begin{align*}
p&:=-\frac{\partial U}{\partial V}=\frac{2}{d}\frac{U}{V}=\frac{2}{d}\frac{2}{(2\pi)^d}\frac{\hbar^2k_F^{d+2}}{2m}\frac{d\tecxt\text{Vol}_d}{d+2}\kom\overset{\substack{{d=2} \\ |}}{=}\frac{\hbar^2k_F^4}{8\pi m}=\frac{U}{V}\\
\kappa_T&=-\frac{1}{V}\frac{\partial V}{\partial p}=-\frac{1}{V}\frac{\partial V}{\frac{\partial U}{V}-\frac{U}{V^2}\partial V}=\frac{1}{-\frac{\partial U}{\partial V}+\frac{U}{V}}=\frac{1}{2p}
\end{align*} 
\end{document}